The following key game mechanics are implemented in the game:

\subsection{Space Movement, Combat and Player Interaction}\label{subsec:movement}

The gameplay features 3 major states - Zero G player movement, Spaceship (entry, flying and shooting), and terrestrial
movement.
This subsection will cover the first two states, with the third being covered in the next subsection.

\subsubsection{Zero G Player Movement and Spaceship Entry}\label{subsubsec:zerog}
The player, when not in their ship, is able to move in 3D world space, with the ability to move in any direction and
rotate in any direction.
This is achieved using Unity's InputSystem\citep{unity:inputsystem} to get the player's input, and then calculating the
torque and thrust force, which are applied to the player's Rigidbody\citep{unity:rigidbody} component.
The player also glides smoothly through space to simulate the lack of friction.
The player also carries a gun, which can be shot using the left mouse button.

From this state, the player can enter their ship by looking at it (when in range) and pressing \textit{F} (the
interact key in the game).
This interaction is performed by repeatedly shooting a ray from the player's camera forwards, and if said ray hits an
interactable object (such as the player's ship) the player is able to interact.
On interaction, an event is triggered.
The interaction event for the ship causes the player to enter the ship and
switch to the spaceship state (wherein the player controls the ship).

The player also has oxygen and health, which are shown in the UI\@.
Oxygen is consumed over time, and regenerates when the player is in the ship or within a planet's atmosphere.

\begin{figure}[h]
    \centering
    \includegraphics[width=0.5\textwidth]{images/player}
    \caption{The player in Zero G movement}
    \label{fig:zerog}
\end{figure}

\subsubsection{Spaceship Movement and Combat}\label{subsubsec:spaceship}

The ship movement is very similar to the player movement, with the player able to move in 3D space, rotate in any
direction, but with a few key differences:
The thrust is much greater in the forwards direction, simulating the fact that the boosters are at the back, and
the pitch and yaw movement is much slower, simulating the fact that the ship is much larger and heavier than the
player, thus requiring more force to rotate.
The player is also able to fire the ship's weapons, which are much more powerful than the player's gun, and have
a cooldown time.
The cooldown timer is shown in the UI\@.

\begin{figure}[h]
    \centering
    \includegraphics[width=0.5\textwidth]{images/ship}
    \caption{The player's ship view}
    \label{fig:ship}
\end{figure}

The player is able to exit the ship by pressing \textit{F} again, allowing the player to return to the Zero G movement
state.

\subsection{Weapon Firing}\label{subsec:weapons}

The player is able to fire their gun and the ship's weapons from their respective states.
These weapons are raycast based, with the player's gun firing a ray from the camera, and the ship's weapons firing
from the front of the ship, with line renderers\citep{unity:linerenderer} being used to simulate the laser beams.
To give a more immersive feel, the player's weapon has a muzzle flash and a sound effect when fired.
If a ray hits an object of a particular layer, the object is damaged, with the damage being calculated based on the
weapon's damage value.
If the objects health reaches 0, the object is destroyed, and can drop resources or other items.
An example of destroyable objects are asteroids, which spawn in a constant rate around the player and are object
pooled for performance reasons (meaning they are reused when destroyed).
Further in development, these asteroids will be crucial for the player to mine resources to repair their ship and
upgrade their equipment.

\begin{figure}[h]
    \centering
    \includegraphics[width=0.5\textwidth]{images/weapon}
    \caption{The player's gun firing}
    \label{fig:weapon}
\end{figure}

\subsection{Resource Collection}\label{subsec:resources}
Resources are instantiated when an object is destroyed, and the player is able to collect these resources by
looking at them and pressing \textit{F} (the interact key).
These resources get added to the player's inventory, which can be accessed by pressing \textit{I}.
The inventory currently shows the player's current resources and the amount of each, as well as sections for the player's
ship and equipment, which will be used for upgrades in the future.
Opening the inventory pauses the game.

\begin{figure}[h]
    \centering
    \includegraphics[width=0.5\textwidth]{images/inventory}
    \caption{The player's inventory}
    \label{fig:inventory}
\end{figure}

\subsection{Planetary Exploration and Basic Life AI}\label{subsec:planetary}

\subsubsection{Planetary Exploration}\label{subsubsec:planetary}

From the ship, the player is able to land on planets or large asteroids, where they can exit the ship and explore
the surface with the effects of gravity.
This requires a change to the player's movement to simulate the gravity, with the addition of
a jump mechanic and more fluid pitch and yaw movement.
A jetpack is also available allowing the player to fly within this gravity affected environment.

As well as being able to land the ship on the planet, if the player is in Zero G movement and gets close enough
to a planet, they will collide with a trigger from the planet's atmosphere.
This will cause the player to be affected by the planet's gravity (and fall to the surface).

In order to simulate the gravity, a custom movement system was implemented to allow for more control over the player's
movement.
This involves applying a force to the player's Rigidbody component based on the player's input, as well
as a rotation force based on the surface normal of the planet.
This ensures the player is always upright with respect to the planet's surface.
Initially, this was only performed when the player is grounded, but was changed to be applied whenever the
player is within the planet's atmosphere, to allow for more fluid movement.
Sounds are played on the player's movement, with footstep sounds being played when the player is walking on the
planet's surface, and jetpack sounds being played when the player is flying.

\begin{figure}[h]
    \centering
    \includegraphics[width=0.5\textwidth]{images/planet}
    \caption{The player on a planet}
    \label{fig:planet}
\end{figure}

\subsubsection{Character Animation}\label{subsubsec:animation}

When the player is walking on a planet, the player's character will animate, with the player's legs and arms moving
in a walking motion.
This is done using Unity's Animator component\citep{unity:animator}, which is called programmatically to change
the player's animation based on the player's movement.
The animations are created using Mixamo\citep{Mixamo:home}, which provides a variety of animations for free.

\subsubsection{Basic Life AI}\label{subsubsec:life}

A basic life AI was implemented to simulate life on the planets.
Aliens wander around the planet's surface, and will follow and stare at the player if they get close enough.
This was done to give an eerie feeling to the planets, as well as to give the player something to interact with.
This is implemented with a state machine approach, with the life AI having a wandering state and a following state,
and various transitions between the states.
An example of one such transition is if the player being close enough to an alien, the alien will switch to the following state.
The alien is also damaged by the player's weapons, and can drop resources when killed.
Again, animations are provided for walking and idle states, similar to the player's character.
Additionally, the alien has a simple pathfinding system to avoid obstacles and follow the player, which involves checking
for obstacles in the direction of the player and moving around them.
Occasionally the alien will emit a sound on interaction, and will also emit footstep sounds when walking.

\begin{figure}[h]
    \centering
    \includegraphics[width=0.5\textwidth]{images/life}
    \caption{The life AI on a planet}
    \label{fig:life}
\end{figure}